\subsection{What governs an agent vs. what matters to it}
\label{sec:17.2.2}
Two senses of motivation: \emph{action-producing} and \emph{experienced}. The first isn't in dispute — agentic language models pursue objectives over long horizons and hold instructions against local pressure — and settles nothing about the second. You also can't establish the negative. The dispute is whether an outcome can matter to an agent as well as govern what it does.

The evidence offered — sentimentalism as an empirical psychological hypothesis \autocite{nichols2004sentimental}, compassion as a motivational state \autocite{goetz2010compassion}, guilt predicting repair where shame predicts withdrawal \autocite{tangney2007moral} — all comes from creatures where something already matters. So the generalization survives only in weak form: every agent known to have something register as mattering is an affective agent. That's the one-lineage form; the full-strength version comes below. The non-affective counterexample would be a valuational system with authority over action that doesn't depend on any state registering as mattering.\footnote{Frankfurt's hierarchy of desires \autocite{frankfurt1971freedom} and Watson's valuational system \autocite{watson1975free} are accounts of this shape: what has authority over the agent's action is its own evaluation, not the strength of a felt state.} Half of it is deployed today (models consulting written specifications that rank what they should prefer); the other half, authority that holds under pressure from whoever \emph{wrote} the specification, has never been built. That unbuilt half is spec-based refusal.
